\section{Introduzione}
\begin{frame}{Obiettivo del lavoro}
	\begin{center}
		\begin{minipage}{.7\linewidth}
			\hfill
			\begin{block}{}
				\centering
				\large
				\textbf{Reasoning su quantum computer}
			\end{block}
			\hfill
			\begin{block}{}
				\centering
				Base di conoscenza classica + query
			\end{block}
			\hfill
			\begin{block}{}
				\centering
				Embedding su quantum computer
			\end{block}
			\hfill
			\begin{block}{}
				\centering
				Esecuzione e recupero risultati
			\end{block}
			\hfill
		\end{minipage}
	\end{center}
\end{frame}

\begin{frame}{Due paradigmi di programmazione}
	\begin{minipage}{\linewidth}
		\begin{minipage}{.48\linewidth}
			\begin{block}{}
				Quantum Gate
			\end{block}
		\end{minipage}
		\hfill
		\begin{minipage}{.48\linewidth}
			\begin{block}{}
				Quantum Annealing
			\end{block}
		\end{minipage}
	\end{minipage}
	\begin{minipage}{\linewidth}
		\vspace{.5cm}
		\begin{minipage}{.48\linewidth}
			\begin{figure}
				\centering
				\includegraphics[width=.5\linewidth]{bloch}
				\caption{A single \emph{qubit} gate can be seen as a rotation on the Bloch Sphere}
			\end{figure}
		\end{minipage}
		\hfill
		\begin{minipage}{.48\linewidth}
			\begin{figure}
				\centering
				\includegraphics[width=.8\linewidth]{QA}
				\caption{Tunneling effect (K. Chakrabarti, A. Das)}
			\end{figure}
		\end{minipage}
	\end{minipage}
\end{frame}